\section{分布式状态提取、恢复与重放}
\label{sec:midterm-state-replay}

一致暂停使各分区停在同一个DUT逻辑周期，但此时的寄存器和存储器状态仍分散在不同FPGA中。要检查这一时刻的完整设计状态，还需要分别读取各分区的数据，并确定每段数据对应哪个逻辑状态。后续恢复与重放也依赖这种对应关系。

REMU已通过扫描链实现内部状态的导出与写回，并支持基于检查点的软件RTL重放\cite{chen_remu_2023}。REMUS复用这些机制，主要完成两方面扩展：一是将划分控制与原有扫描链控制连接起来，使扫描操作能够在多FPGA一致暂停后进行；二是记录面向分区的状态映射，将各FPGA导出的数据组织为可供恢复和重放使用的全局检查点。

\subsection{划分控制与扫描链控制的协同}
\label{sec:midterm-scan-control}

\paragraph{将跨分区传播接入运行控制}

扫描链控制器除组织扫描操作外，还通过运行、暂停和扫描模式控制DUT的执行。设计划分后，这部分运行控制需要知道跨片传播是否完成。例如，某分区的本地模型已经完成交互，但其他FPGA上的组合计算结果尚未到达，此时仍不能推进DUT逻辑时钟。若仅按本地完成条件更新周期计数，计数值也会与DUT实际完成的周期数不符。

REMUS将划分控制机产生的逻辑周期完成信号接入扫描链侧的运行控制。本地模型交互完成后，划分控制机按LI-TDM流程组织跨片传播；最后一个传播阶段完成时，才允许DUT更新时序状态。运行控制据此增加逻辑周期计数，并递减剩余执行步数。因此，DUT状态更新、周期计数和定步执行使用同一个周期完成条件。

两个控制部分还交换运行状态和暂停协调信息。运行控制向划分控制机提供运行使能及同步屏障请求；划分控制机返回当前传播阶段和屏障确认。暂停请求到达时，如果分区尚处于周期内的传播过程，运行控制就等待该周期完成后再停止正常执行。主机则通过屏障确认和周期计数协调各分区。第\ref{sec:midterm-global-pause}节所述的一致暂停过程建立在这些连接之上。

\paragraph{从一致暂停转入扫描}

扫描开始前，主机先完成各分区的一致暂停，并等待相关模型进入空闲状态。随后，主机将各分区切换到扫描模式，启动局部扫描控制。这里保留了REMU已有的扫描时序，新增的协调工作是使所有分区在同一逻辑周期停止正常执行后，再进行各自的扫描操作。

正常执行和扫描访问采用不同的时钟使能条件。正常执行时，DUT状态只在逻辑周期完成条件满足后更新；扫描时，则由扫描使能驱动状态移位和存储器访问。扫描会消耗FPGA物理时钟周期，但不增加DUT逻辑周期计数。主机等待全部分区的扫描操作结束后退出扫描模式，再根据后续命令保持暂停或继续运行。

\subsection{分布式状态提取与全局检查点组织}
\label{sec:midterm-global-snapshot}

\paragraph{局部扫描数据的存储与收集}

REMUS在每个分区内保留独立的扫描结构。扫描导出（Scan Dump）时，各分区的数据先写入各自FPGA上的内存，主机随后按分区读取并保存。这样，多个FPGA不需要串接成一条跨片扫描链，各分区可以分别完成本地状态访问。

以双FPGA设计为例，FPGA A和FPGA B均在逻辑周期$T$完成一致暂停。A先完成扫描导出、B稍后完成，只要期间两者均未恢复正常执行，导出的状态就仍对应同一逻辑时刻。主机可以依次读取两份数据，而不必要求它们在同一个物理时钟周期到达。这里的前提是第\ref{sec:midterm-global-pause}节所述的有效暂停条件已经满足，扫描操作也保持REMU原有的状态导出行为。

主机将同一逻辑周期的各分区扫描数据与模型内存等状态共同组织为全局检查点（Global Snapshot）。各分区的扫描数据保存在独立文件中，并记录所属分区和检查点周期。全局检查点表示完整设计在该周期的状态集合，无须将全部数据拼接成一个文件。

\paragraph{从局部扫描位置映射到逻辑状态}

只有扫描数据还不足以解释设计状态。扫描文件中的一段比特本身不包含寄存器名称，也不说明它属于哪个分区。例如，FPGA A与FPGA B的扫描数据都可能从文件起始位置保存一个寄存器；如果缺少分区归属和局部布局，软件便无法确定这两个值在完整设计中各自对应的对象。

为此，REMUS在编译时分别记录每个分区的状态布局，再将它们合并为全局状态映射元数据。分区标识用于选择局部数据，扫描顺序、位宽和片段偏移用于确定取出哪些比特，层次名称则用于找到对应的逻辑状态。对于存储器，还需记录其字宽、深度等信息，以还原各存储单元的内容。

解析时，软件先读取指定分区的扫描数据，再按布局取出各状态片段，将其放入对应的寄存器或存储器对象中。经过这一处理，分散在各FPGA中的扫描数据就可以按完整设计的状态名称统一访问。状态名称与编译生成的重放RTL保持对应，包括划分后形成的层次结构，因此检查点、状态元数据和重放RTL需要配套使用。

\paragraph{检查点组成与既有功能复用}

除扫描数据外，后续执行还需要存储模型中的程序和数据。例如，处理器寄存器已恢复，但程序所在的模型内存未恢复，处理器仍无法从该检查点正确继续执行。因此，保存检查点时还需按分区收集相应模型内存。对于后续发生变化的外部输入，则沿用REMU已有的输入记录机制，按逻辑周期重现这些变化。表\ref{tab:midterm-checkpoint-content}列出了检查点及其配套信息。

\begin{table}[htbp]
    \centering
    \caption{分布式检查点及其配套信息}
    \label{tab:midterm-checkpoint-content}
    \small
    \renewcommand{\arraystretch}{1.25}
    \begin{tabularx}{\linewidth}{@{}>{\raggedright\arraybackslash}p{0.23\linewidth}>{\raggedright\arraybackslash}X@{}}
        \toprule
        内容      & 用途                                          \\
        \midrule
        分区扫描数据  & 保存局部寄存器、RTL存储器及扫描覆盖的模型状态，先写入各自FPGA内存，再由主机收集 \\
        模型内存数据  & 保存程序、数据等存储模型内容，供恢复后继续执行使用                   \\
        状态映射元数据 & 将分区内的扫描位置对应到逻辑状态，支持统一解码和按分区写回               \\
        周期与区域信息 & 记录检查点周期、分区数量及内存区域归属，定位待加载的数据                \\
        输入记录    & 保存受管理输入的取值和变化周期，重现检查点之后的输入变化                \\
        \bottomrule
    \end{tabularx}
\end{table}

在完成控制协调和分区状态映射后，REMUS即可将这些数据接入REMU已有的恢复与重放流程。向FPGA恢复时，软件按分区归属写回数据，并通过扫描写入恢复DUT状态；在软件RTL中重放时，软件将各分区的扫描数据解码为逻辑状态，再加载到对应仿真对象中。两条路径分别复用原有的Scanin和软件重放机制，软件重放也继续使用原有的波形观测方式。本节的扩展重点是使这些机制能够正确定位和使用多FPGA上的状态。

\subsection{功能验证与阶段性结果}
\label{sec:midterm-state-validation}

\paragraph{扫描状态与单体参考的对照}

状态提取实验需要检查两个问题：各分区是否导出了共同暂停周期的状态，以及软件是否正确解释了这些数据。实验以\allowbreak\mbox{PicoRV32}\allowbreak{}和香山昆明湖（KMH）为DUT，分别运行划分后的Cosim与未划分的单体RTL仿真。多分区设计完成一致暂停后，通过Scan Dump导出内部状态，将数据落盘并按分区元数据解码；单体仿真执行至相同逻辑周期，提供对应状态作为参考。

已有实验中，两种设计在已比较的内部状态范围内结果相同。这一比较经过了扫描导出、文件保存和状态解码过程，说明所得结果覆盖了从局部扫描到逻辑状态解析的完整路径。它与第\ref{sec:midterm-cycle-cosim}节的暂停验证使用相同的周期基准，本节进一步确认了用于比较的状态能够通过分布式扫描获取。

% TODO(状态提取实验)：补充各DUT的版本、程序、分区配置、暂停周期，
% 以及实际比较的寄存器/存储状态范围、比较方式和日志来源。
% 上述定性结果及完整Scan Dump、落盘、解码链路已经作者确认；
% 不将已比较的状态范围扩写为所有内部对象均已逐位验证。

\paragraph{非零周期检查点的恢复与重放}

为检查分区检查点能否被既有恢复与重放流程使用，实验采用运行过程中保存的非零周期检查点，分别在逻辑仿真软件中进行Replay，并在FPGA上进行Replay和Restore。验证关注加载后继续执行的结果，而不只检查数据能否成功写入。

已有测试中，软件Replay及FPGA Replay、Restore后的运行结果均符合预期。继续运行后再次保存检查点，其中的数据也符合预期。这些结果表明，在已测试的条件下，分区状态能够接入已有恢复与重放流程，并用于后续执行。

% TODO(恢复与重放实验)：补充每组实验的DUT、分区数量、仿真器或板卡，
% 起始检查点周期、后续执行区间、预期结果来源与再次保存的数据比较范围。
% 非零周期Replay/Restore及后续检查点数据符合预期已经作者确认；
% 不直接认定这些实验均覆盖PicoRV32和KMH，也不填入未经提供的性能数据，
% 不将运行结果符合预期扩写为全部内部信号的波形已经逐周期对照。
